\documentclass[10pt]{article}
\usepackage[T1]{fontenc}
\usepackage[utf8]{inputenc}
\usepackage[french]{babel}
\usepackage{amsmath,amssymb,mathrsfs,mathpazo}
\usepackage{geometry}
\geometry{a4paper,margin=2cm}
\usepackage{graphicx,array,tabularx,tabularray,colortbl,multirow,multicol,booktabs,enumitem,pifont,ulem,fancybox,diagbox,cancel}
\usepackage{pstricks,pstricks-add,pst-plot,pst-tree,pst-node,pst-text,pst-func,pst-math,pst-bspline,pst-3dplot,pst-eucl,pst-solides3d}
\let\pscalculate\undefined
\usepackage{pst-bezier}
\usepackage{tikz,pgfplots,tkz-tab}
\usetikzlibrary{arrows,arrows.meta,patterns,calc,positioning,trees,decorations.markings,shapes.geometric,shapes.misc,shapes.arrows,intersections,decorations.pathmorphing,decorations.pathreplacing,backgrounds,shadows,babel}
\pgfplotsset{compat=1.18}
\usepgfplotslibrary{fillbetween}
\usepackage[np]{numprint}
\usepackage{listings,makecell,pas-tableur,twoopt}
\let\barre\undefined
\usepackage[pstricks]{tablvar}
\definecolor{darkgreen}{RGB}{0,100,0}
\providecommand{\red}{\color{red}}
\providecommand{\blue}{\color{blue}}
\providecommand{\green}{\color{green}}
\providecommand{\black}{\color{black}}
\providecommand{\white}{\color{white}}
\providecommand{\gray}{\color{gray}}
\providecommand{\cyan}{\color{cyan}}
\providecommand{\magenta}{\color{magenta}}
\providecommand{\yellow}{\color{yellow}}
\providecommand{\violet}{\color{violet}}
\providecommand{\orange}{\color{orange}}
\providecommand{\euro}{\texteuro}
\providecommand{\EUR}{\texteuro}
\providecommand{\ieme}{\textsuperscript{e}}
\providecommand{\degres}{\ensuremath{^\circ}}
\providecommand{\degree}{\ensuremath{^\circ}}
\providecommand{\vv}[1]{\overrightarrow{#1}}
\providecommand{\Coord}[1]{#1}
\providecommand{\decosix}{}
\providecommand{\starredbullet}{\ensuremath{\bullet}}
\pagestyle{empty}
\setlength{\parindent}{0pt}
\usepackage{adjustbox}\DeclareUnicodeCharacter{2236}{:}\DeclareUnicodeCharacter{2190}{\ensuremath{\leftarrow}}
\providecommand{\R}{\mathbb{R}}
\providecommand{\N}{\mathbb{N}}
\providecommand{\Z}{\mathbb{Z}}
\providecommand{\Q}{\mathbb{Q}}
\providecommand{\D}{\mathbb{D}}
\providecommand{\C}{\mathbb{C}}
\providecommand{\e}{{\mathrm{e}}}
\providecommand{\vect}[1]{\overrightarrow{#1}}
\providecommand{\vectt}[1]{\overrightarrow{\mathrm{#1}}}
\providecommand{\vvt}[1]{\overrightarrow{\mathrm{#1}}}
\providecommand{\np}[2][]{{#2~#1}}
\providecommand{\vv}[1]{{\overrightarrow{#1}}}
\renewcommand{\d}{{\mathrm{d}}}
\providecommand{\degres}{{^\circ}}
\providecommand{\degre}{{^\circ}}
\providecommand{\blue}{\color{blue}}
\providecommand{\red}{\color{red}}
\providecommand{\enskip}{\quad}
\providecommand{\llbracket}{\lbrack\!\lbrack}
\providecommand{\rrbracket}{\rbrack\!\rbrack}
\providecommand{\euro}{{\text{€}}}
\providecommand{\no}{{\text{n°}}}
\providecommand{\interval}[2]{{[#1;#2]}}
\providecommand{\intervFO}[2]{{[#1;#2[}}
\providecommand{\dotfill}{\ldots}

\providecommand{\ds}{\displaystyle}
\providecommand{\textding}[1]{#1}
\providecommand{\Oij}{\(\mathrm{O};\vec{\imath},\vec{\jmath}\)}
\providecommand{\Oijk}{\(\mathrm{O};\vec{\imath},\vec{\jmath},\vec{k}\)}
\providecommand{\Ouv}{\(\mathrm{O};\vec{u},\vec{v}\)}
\providecommand{\intervFF}[2]{\ensuremath{[#1;#2]}}
\renewcommand{\euro}{\ifmmode\text{\texteuro}\else\texteuro\fi}
\DeclareUnicodeCharacter{00B0}{\ifmmode^{\circ}\else\textdegree\fi}
\DeclareUnicodeCharacter{20AC}{\ifmmode\text{\texteuro}\else\texteuro\fi}
% Énoncés : Éducation nationale. Transcription et corrigés : APMEP, Denis Vergès et contributeurs. Découpage et classement : CoopMaths, Sébastien Lozano, Éric Elter et contributeurs. Adaptation de la présentation : Maths Explorer.
% CC BY-SA : https://coopmaths.fr/www/about/
% https://coopmaths.fr/alea/static/e3c/2020/tex/e3c_2020_00_sujet4_1_cor.tex
\begin{document}
\noindent\textit{APMEP / CoopMaths. CC BY-SA. Présentation adaptée par Maths Explorer.}\par\medskip
\section*{Exercice 1 (5 points)}
	
	\subsection*{Question 1}
	Le milieu $I$ de $[AB]$ a pour coordonnées $(3, 4)$. Si $(\Delta)$ est la médiatrice de $[AB]$, alors
	\[
	M(x, y) \in (\Delta) \iff \overrightarrow{IM} \cdot \overrightarrow{AB} = 0 \iff -2(x - 3) + 4(y - 4) = 0 \iff -2x + 4y - 10 = 0 \iff x - 2y + 5 = 0
	\]
	La réponse correcte est \textbf{b.}
	
	\subsection*{Question 2}
	Les vecteurs $\overrightarrow{MP}$ et $\overrightarrow{MN}$ sont orthogonaux, donc les droites $(MP)$ et $(MN)$ sont perpendiculaires : le triangle $MNP$ est donc rectangle en $M$ et l’ensemble des points est le cercle de diamètre $[NP]$.\\ La réponse correcte est \textbf{b.}
	
	\subsection*{Question 3}
	On a $g'(x) = 3x^2 - 4$ et en particulier $g'(-1) = 3 - 4 = -1$.\\ Si $t$ est la tangente, $M(x, y) \in T \iff y - g(-1) = g'(-1)(x - (-1))$.\\ Avec $g(-1) = -1 + 4 + 5 = 8$, on a donc :
	\[
	M(x, y) \in T \iff y - 8 = -1(x + 1) \iff y = -x + 7
	\]
	La réponse correcte est \textbf{c.}
	
	\subsection*{Question 4}
	En écrivant $y = x^2 + x + 3 = \left( x + \dfrac{1}{2} \right)^2 - \dfrac{1}{4} + 3 = \left( x + \dfrac{1}{2} \right)^2 + \dfrac{11}{4}$, on voit que la droite d’équation $x = -0,5$ est axe de symétrie de la parabole.\\
	 La réponse correcte est \textbf{d.}
	
	\subsection*{Question 5}
	L’inéquation peut s’écrire $3e^{x+2} < 3e^4$ ou en simplifiant par 3 :\\ $e^{x+2} < e^4$ et en multipliant par $e^{-4}$, $e^{0} < e^{x-2}$, soit finalement par croissance de la fonction exponentielle, $0 < x - 2$ ou $x < 2$. L’ensemble des solutions est l’intervalle $] -\infty, 2[$.\\
	 La réponse correcte est \textbf{c.}

\end{document}
